\section{Scope and sources}
\label{app:scope}

\subsection{Selection and scope boundaries}
Branch selection combines a measurement question with identifiable historical evidence. UMA and Tellor cover contestable reporting and adjudication; Chainlink covers staking and selected feed service; Flare covers entitlements and claims; Pyth covers publisher--epoch incentives and staking. Chronicle, RedStone and DIA add narrower challenge, Push-update and withdrawal observations. Table~\ref{tab:protocol_sources} defines these boundaries. The registry is not exhaustive or probabilistic; exclusions apply only to this release.

Economic comparison requires an interpretable native unit, identifiable implementation and source. Monetary amounts also require an asset and accounting stage. Signed updates can describe delivery without measuring publisher payments; withdrawals can separate principal and reward without proving price accuracy. These limits govern the research opportunities in Table~\ref{tab:research_opportunities}.

\subsection{Historical implementations and source register}
\label{app:implementations}
Table~\ref{tab:implementation_register} records the principal contracts, programs and source versions. The accompanying source register describes their network roles, observation units, historical boundaries and inclusion criteria, including the Tellor modules and Chainlink feed proxy. Acquisition records preserve the connection between each sampled implementation and the evidence used to reconstruct it~\cite{atlasScopeMetadata2026}.

Deployment, scan and observation boundaries remain separate. Chronicle and RedStone searches begin at block 0, a search bound rather than a deployment block. Pyth's archive starts at a transaction timestamp rather than an EVM block. A pinned repository or cutoff runtime does not prove an unchanged historical implementation. Deployment times and implementation intervals require supporting metadata.

\begingroup
\small\setlength{\tabcolsep}{4pt}\renewcommand{\arraystretch}{1.12}
\begin{longtable}{@{}L{0.12\linewidth} L{0.40\linewidth} L{0.42\linewidth}@{}}
\caption{Historical implementation and source register. Declared deployment blocks and retained source versions reproduce the deposited metadata; they are not newly verified deployment or runtime histories. Full addresses, source commits and unresolved fields are preserved in the accompanying files.}\label{tab:implementation_register}\\
\toprule
\textbf{Branch} & \textbf{Principal identifier and chain role} & \textbf{Retained source, boundary and limitation} \\
\midrule\endfirsthead
\multicolumn{3}{l}{\small\itshape Table \thetable\ continued}\\
\toprule
\textbf{Branch} & \textbf{Principal identifier and chain role} & \textbf{Retained source, boundary and limitation} \\
\midrule\endhead
\midrule\multicolumn{3}{r}{\small\itshape Continued on next page}\\\endfoot
\bottomrule\endlastfoot
\textbf{UMA} & Polygon (137), OOv2: \path{0xee3afe347d5c74317041e2618c49534daf887c24}. Ethereum (1), VotingV2: \path{0x004395edb43efca9885cedad51ec9faf93bd34ac}. Registered adapters and tunnels are listed separately in the contract CSV. Roles: verification, settlement and delivery, by implementation. & Retained UMA source \texttt{a16ee53125c4} and adapter source \texttt{8b76cc9e0d46}. Declared OOv2 deployment: Polygon block 29,786,052; earliest registered adapter: 23,569,780; cutoff: 89,439,336. DVM: 16,697,232--25,433,938. Historical runtime intervals are not fully established by cutoff code hashes. \\ \addlinespace[3pt]
\textbf{Chainlink} & Ethereum (1), community pool: \path{0xbc10f2e862ed4502144c7d632a3459f49dfcdb5e}; operator pool: \path{0xa1d76a7ca72128541e9fcacafbda3a92ef94fdc5}; ETH/USD proxy: \path{0x5f4ec3df9cbd43714fe2740f5e3616155c5b8419}. Roles: staking accounting and feed delivery. & Staking v0.2 application configuration and ABI; retained phase-specific feed evidence. Pools start at declared deployment block 18,572,190; vault at 18,621,067; controller at 18,621,068; cutoff 25,433,938. Exact historical frontend/code version bindings remain incomplete. \\ \addlinespace[3pt]
\textbf{Tellor} & \texttt{tellor-1} application modules, ABCI \texttt{new\_report} events and block results. Roles: reporting, adjudication and economic accounting on Tellor Layer. & Retained Layer source \texttt{943a2709ef0a}. Registered heights 1--19,890,860; first report at 317. Public opening is not genesis. The pre-opening network stage and completeness across historical module changes require separate evidence. \\ \addlinespace[3pt]
\textbf{Flare FTSOv2} & Flare (14), RewardManager: \path{0xc8f55c5aa2c752ee285bd872855c749f4ee6239b}; VoterRegistry: \path{0x2580101692366e2f331e891180d9ffdf861fce83}. Roles: provider/registry evidence and reward accounting. & Reward artifacts: \texttt{fsp-rewards} commit \texttt{56dccda13597}. Declared deployment blocks 29,549,020 and 29,549,006, respectively; cutoff 64,054,789. Closed epochs 228--410; provider conditions start at 251. Aggregate claims need not identify every component. \\ \addlinespace[3pt]
\textbf{Pyth} & Solana mainnet OIS program: \path{pyti8TM4zRVBjmarcgAPmTNNAXYKJv7WVHrkrm6woLN}. Role: staking and transfer accounting; Pythnet-derived states are retained inputs, not a full independent chain scan. & Retained governance source \texttt{68a9a36ec3d4}. Transaction timestamps: 2024-09-11 19:53:53 to 2026-06-30 23:11:29 UTC. State epochs 2898--2946 cover 49 closed epochs. A later snapshot does not by itself establish historical decision-time availability. \\ \addlinespace[3pt]
\textbf{Chronicle} & Ethereum (1): 72 validated Scribe/ScribeOptimistic emitted-event contracts, individually identified in the source register. Roles: report validation, delivery and challenge-reward accounting. & Source \texttt{12ff06ca7881}; validator registry \texttt{f06e0a8209fc}. Registered scan: blocks 0--25,433,938. Zero is a scan start, not deployment. Historical deployment dates and per-interval implementation bindings are not established for every address. \\ \addlinespace[3pt]
\textbf{RedStone Push} & Ethereum (1): 9 Push adapters and 91 feed contracts in the frozen registry, with individual addresses recorded in the source register. Roles: signed-publication verification and delivery. & Source/manifest commit \texttt{95f8cb9fd00b}. Registered scan: blocks 0--25,433,938; deployment dates are separate. First retained update is not first protocol launch. Pull modes, other networks and full AVS economics are outside this branch. \\ \addlinespace[3pt]
\textbf{DIA} & Lasernet (1050), staking: \path{0x677Cf1299c367F6cf6F3E1669aCC18Fd059a5919}; wrapped-DIA address is retained in the contract CSV. Role: staking accounting and principal/reward payment. & Staking-application configuration and ABI retained in the manifest. Declared staking deployment: block 3,874,186; cutoff: 31,302,665. Deployment UTC time and a complete historical runtime-version mapping are not established by this summary. \\ \addlinespace[3pt]
\end{longtable}
\endgroup


\subsection{Evidence for the ecosystem figure}
\label{app:ecosystem_evidence}
Figure~\ref{fig:ecosystem_scope} uses first and last timestamps from the tables named in Table~\ref{tab:figure_observation_windows}, derived from pinned coverage metadata~\cite{atlasScopeMetadata2026}. These extents do not establish full branch histories or continuous coverage. Pyth uses epoch-end labels, while its transaction archive starts earlier. Tellor reports predate the public-opening announcement. Chainlink feed and DIA withdrawal histories need not start at staking deployment or launch.

\begingroup
\small\setlength{\tabcolsep}{4pt}\renewcommand{\arraystretch}{1.12}
\begin{longtable}{@{}L{0.17\linewidth} L{0.40\linewidth} L{0.18\linewidth} L{0.18\linewidth}@{}}
\caption{Selected-table timestamp extents used in Figure~\ref{fig:ecosystem_scope}. Dates are UTC calendar dates, not deployment or scan bounds. Exact UTC timestamps and table identifiers are supplied with the figure source data.}\label{tab:figure_observation_windows}\\
\toprule
\textbf{Branch} & \textbf{Selected observation} & \textbf{First retained date} & \textbf{Last retained date} \\
\midrule\endfirsthead
\multicolumn{4}{l}{\small\itshape Table \thetable\ continued}\\
\toprule
\textbf{Branch} & \textbf{Selected observation} & \textbf{First retained date} & \textbf{Last retained date} \\
\midrule\endhead
\midrule\multicolumn{4}{r}{\small\itshape Continued on next page}\\\endfoot
\bottomrule\endlastfoot
UMA & Request rounds & 2022-11-04 & 2026-06-30 \\ \addlinespace[3pt]
Chainlink & ETH/USD service windows & 2023-11-14 & 2026-06-30 \\ \addlinespace[3pt]
Tellor & Micro-reports & 2025-05-09 & 2026-06-30 \\ \addlinespace[3pt]
Flare FTSOv2 & Reward-claim events & 2024-10-03 & 2026-06-30 \\ \addlinespace[3pt]
Pyth & Publisher epoch-end labels & 2025-07-24 & 2026-06-25 \\ \addlinespace[3pt]
Chronicle & Validator / challenge events & 2023-08-28 & 2026-06-30 \\ \addlinespace[3pt]
RedStone Push & Push update events & 2023-12-07 & 2026-06-30 \\ \addlinespace[3pt]
DIA & Staking withdrawals & 2025-07-08 & 2026-06-18 \\ \addlinespace[3pt]
\end{longtable}
\endgroup


Table~\ref{tab:historical_milestones} identifies dated historical markers at their source precision. Announcements, public access and deployment remain distinct: UMA's optimistic-oracle launch, for example, does not date OOv2 deployment. Each marker requires a source identifying the event and its date.

\begingroup
\small\setlength{\tabcolsep}{4pt}\renewcommand{\arraystretch}{1.12}
\begin{longtable}{@{}L{0.19\linewidth} L{0.18\linewidth} L{0.58\linewidth}@{}}
\caption{Historical markers in Figure~\ref{fig:ecosystem_scope}, dated at source precision. Announcements do not establish deployment. No verified RedStone first-Push milestone is available for inclusion.}\label{tab:historical_milestones}\\
\toprule
\textbf{Branch} & \textbf{Date / precision} & \textbf{Event and primary source} \\
\midrule\endfirsthead
\multicolumn{3}{l}{\small\itshape Table \thetable\ continued}\\
\toprule
\textbf{Branch} & \textbf{Date / precision} & \textbf{Event and primary source} \\
\midrule\endhead
\midrule\multicolumn{3}{r}{\small\itshape Continued on next page}\\\endfoot
\bottomrule\endlastfoot
UMA & 2018-12 & Whitepaper; Project / network history~\cite{risklabsHistory2026}. \\ \addlinespace[3pt]
UMA & 2021-05-13 & Optimistic Oracle; Service / public access~\cite{umaOOlaunch2021}. \\ \addlinespace[3pt]
Chainlink & 2019-05-30 & Ethereum mainnet; Service / public access~\cite{chainlinkMainnet2019}. \\ \addlinespace[3pt]
Chainlink & 2023-11-28 & Staking v0.2; Mechanism announcement / access~\cite{chainlinkV022023}. \\ \addlinespace[3pt]
Tellor & 2019-08-05 & Ethereum mainnet; Service / public access~\cite{tellorMainnet2019}. \\ \addlinespace[3pt]
Tellor & 2025-08-04 & Layer public access; Mechanism announcement / access~\cite{tellorPublic2025}. \\ \addlinespace[3pt]
Flare FTSOv2 & 2022-07-14 & Flare genesis; Project / network history~\cite{flareGenesis2025}. \\ \addlinespace[3pt]
Flare FTSOv2 & 2024-09-16 & FTSOv2 announcement; Mechanism announcement / access~\cite{flareV22024}. \\ \addlinespace[3pt]
Pyth & 2021-08-26 & Solana feeds; Service / public access~\cite{pythMainnet2021}. \\ \addlinespace[3pt]
Pyth & 2024-09-02 & OIS introduction; Mechanism announcement / access~\cite{pythOISintro2024}. \\ \addlinespace[3pt]
Chronicle & 2017 & Maker oracle roots; Project / network history~\cite{chroniclePublic2023}. \\ \addlinespace[3pt]
Chronicle & 2023-09-05 & Public access; Service / public access~\cite{chroniclePublic2023}. \\ \addlinespace[3pt]
DIA & 2025-03-26 & Lasernet mainnet; Project / network history~\cite{diaMainnet2025}. \\ \addlinespace[3pt]
DIA & 2025-06-24 & Staking launch; Mechanism announcement / access~\cite{diaStaking2025}. \\ \addlinespace[3pt]
\end{longtable}
\endgroup


The chain matrix distinguishes the sampled mechanisms from the wider services described in protocol documentation~\cite{uma2026,chainlinkstaking2026,tellorDisputes2026,flareFTSO2026,pythOIS2026,chronicleDocs2026,redstoneDocs2026,diaDocs2026}. Its historical scope remains defined by the sampled implementations. Pythnet identifies retained inputs used in the Pyth reconstruction; it does not denote a complete chain scan. Unmarked cells indicate relationships outside the documented figure scope.
